\section{Conclusion}
\label{sec:conclusion}

Power-law lag geometry under signal-coupled heteroscedastic noise is governed by
one information scale, $n\sum_{j \leq L} j^{2H}$. That quantity controls
estimation, the residual adequacy law, and the minimax separation boundary, with
fixed-$L$ and growing-$L$ regimes as two manifestations of the same geometry.

The simulation evidence supports the theory at the levels of calibration,
oracle agreement, and separation-bound behavior. The information-normalized
error constant is approximately $2.22$ on the tested benchmark.

For unknown $\sigma_0$, the decoupled benchmark is an exact split-$F$ law,
while the feasible same-sample extension is operational under parametric
bootstrap calibration.

The joint asymptotic law of $(\hat\alpha,\Hhat)$ also makes the lag geometry
available to downstream smooth functionals, including plug-in memory horizons
studied separately.
